\documentclass[11pt]{amsart}
\usepackage[localtheorems]{raeez-math-template}

\providecommand{\C}{\mathbb{C}}
\providecommand{\R}{\mathbb{R}}
\providecommand{\Z}{\mathbb{Z}}
\providecommand{\cA}{\mathcal{A}}
\providecommand{\cE}{\mathcal{E}}
\providecommand{\cO}{\mathcal{O}}
\providecommand{\cD}{\mathcal{D}}
\providecommand{\cF}{\mathcal{F}}
\providecommand{\cG}{\mathcal{G}}
\providecommand{\cJ}{\mathcal{J}}
\providecommand{\Obs}{\mathrm{Obs}}
\providecommand{\dbar}{\bar{\partial}}
\providecommand{\g}{\mathfrak{g}}
\providecommand{\tr}{\mathrm{tr}}
\providecommand{\gh}{\mathrm{gh}}
\providecommand{\BV}{\mathrm{BV}}
\providecommand{\hCS}{\mathrm{hCS}}
\providecommand{\kappach}{\kappa_{\mathrm{ch}}}
\providecommand{\kappacat}{\kappa_{\mathrm{cat}}}
\providecommand{\Map}{\mathrm{Map}}
\providecommand{\PTVV}{\mathrm{PTVV}}
\providecommand{\Crit}{\mathrm{Crit}}

\theoremstyle{plain}
\newtheorem{theorem}{Theorem}[section]
\newtheorem{proposition}[theorem]{Proposition}
\newtheorem{lemma}[theorem]{Lemma}
\theoremstyle{definition}
\newtheorem{definition}[theorem]{Definition}
\newtheorem{construction}[theorem]{Construction}

\title{4D holomorphic Chern--Simons on a CY$_2$: the $(-2)$-shifted
symplectic partner of 6D hCS}
\author{Wave-15 L4/20}
\date{2026-04-22}

\begin{document}
\maketitle

Wave~14 I4 fixed the CY-dimension law that controls BV parity: on
$\mathrm{Map}(Y_d, BG)$ the Pantev--To\"en--Vaqui\'e--Vezzosi canonical
form is $n$-shifted symplectic with $n = d-4$. At $d=3$ one recovers
the $(-1)$-shift native to hCS on $\C^3$; the present note executes
the $d=2$ reduction, for which the shift is $-2$. The object is
4D holomorphic Chern--Simons on a CY$_2$ surface $X$, and the
construction is not a curiosity: the $(-2)$-shift matches the
Class M $E_2$-chiral stratum of Vol~III at $d=2$ and identifies the
BV bracket with the cobracket on $\kappach$-weight~$2$ observables.

\section{The 4D hCS action on a CY$_2$}

Fix $X$ a compact CY$_2$ (K3, a complex torus $T^4$, or a bielliptic
surface) with nowhere-vanishing holomorphic 2-form $\Omega_2 \in
H^0(X,K_X) = H^{2,0}(X)$, and $\g$ reductive with invariant $\tr$.
The fields are
\[
\cA \;\in\; \cE \;=\; \Omega^{0,\bullet}(X,\g)[1],\qquad
\cA \;=\; c + A_{0,1} + A^*_{0,2} \quad\text{(ghosts and antifields
encoded in Dolbeault form-degree).}
\]
The classical action is
\[
S_\hCS^{(4)}(\cA) \;=\; \int_X \Omega_2 \wedge \tr\!\Bigl(
 \tfrac12\, \cA \wedge \dbar\cA \;+\; \tfrac13\, \cA\wedge\cA\wedge\cA
\Bigr),
\]
a $(2,2)$-form integrated against the CY$_2$ volume. The Lie triple
bracket $[\cA,\cA,\cA]_{\mathrm{Lie}_3}$ reduces to the ordinary
wedge-Lie cube because $\g$ is associative on the matrix side; the
subscript is retained only to mark that the cyclic invariance
$\tr[a,b,c] = \tr[b,c,a]$ is what licences the symmetric third-order
term. Variation yields
\[
\mathrm{EL}(S_\hCS^{(4)}) \;=\; \dbar A_{0,1} + A_{0,1}\wedge A_{0,1}
\;\in\; \Omega^{0,1}(X,\g),
\]
the flatness equation of a partial holomorphic connection, whose
critical locus is $\Map_{\mathrm{hol}}(X_{\mathrm{dR}}, BG)$ in the
$(0,1)$-gauge.

\section{$(-2)$-shifted symplectic structure}

\begin{theorem}[PTVV shift at $d=2$]
\label{thm:minus2-shift}
The BV pairing
\[
\omega_\BV^{(2)} \;=\; \int_X \Omega_2 \wedge \tr\!\bigl(\delta\cA
 \wedge \delta\cA\bigr)
\]
is a $(-2)$-shifted symplectic form on $\cE$. Equivalently, under the
identification $\Crit(S_\hCS^{(4)}) = \Map(X_{\mathrm{dR}}, BG)$, the
pullback of the PTVV canonical form of \cite{PTVV2013}
on $\Map(X, BG)$ for $X$ a $2$-dimensional Calabi--Yau is
$(-2)$-shifted.
\end{theorem}

\begin{proof}[Sketch]
Serre duality on $X$ pairs $\Omega^{0,k}(X,\g)$ with
$\Omega^{0,2-k}(X,\g^\vee)$ against $\Omega_2$. The quartet collapses
to two layers: $\delta A_{0,1}$ pairs with $\delta A^*_{0,1}$ via
$\Omega^{0,1}\otimes\Omega^{0,1}\to\Omega^{0,2}
\xrightarrow{\int\Omega_2\wedge} \C$; $\delta c$ pairs with
$\delta c^*_{0,2}$ via $\Omega^{0,0}\otimes\Omega^{0,2}\to
\Omega^{0,2}$. Both pairings land in ghost-number $\gh = -2$, not
$-1$: at $d=2$ the ghost $c$ and its antifield $c^*_{0,2}$ live in
Dolbeault degrees $0$ and $2$, summing to $2$, so $\gh(c) + \gh(c^*)
= 1 + (-3) = -2$ under the I4 convention. Equivalently
$\deg \omega_\BV^{(2)} = d - 4 = -2$. Calaque--Pantev--To\"en--
Vaqui\'e--Vezzosi \cite{CPTVV2017} identify this as the
AKSZ transgression of the invariant pairing on $B\g$ along
the CY$_2$ cycle $[X]\in H_{2,2}(X)$.
\end{proof}

\section{$\Obs_{\hCS}^{(4)}$ is an $E_2$-algebra}

\begin{theorem}[$E_2$-observables]
\label{thm:E2-obs}
The factorisation algebra $\Obs_{\hCS}^{(4)}$ of classical observables
of $S_\hCS^{(4)}$ on $X$ is locally constant along the transverse
$\R^2 = \C$-direction inside a Dolbeault-framed tubular neighbourhood,
and carries an $E_2$-algebra structure on its global sections. In the
Costello--Gwilliam formalism \cite{CG2021}
\[
\Obs_{\hCS}^{(4)}(U) \;=\; \bigl(\cO(\cE(U))[[\hbar]],\,
 d_\BV + \hbar \Delta_\BV\bigr),
\qquad U \subset X \text{ open,}
\]
is an $E_2$-factorisation algebra, not $E_3$, because the two
holomorphic directions of $X$ provide only two topological
disc-directions after Dolbeault cohomology is taken.
\end{theorem}

\begin{proof}[Sketch]
Run the Costello--Gwilliam argument of \cite[\S4.8]{CG2021} with
the CY dimension reduced from $3$ to $2$. Each holomorphic direction
contributes one $E_1$-factor after passing to $\dbar$-cohomology; the
two holomorphic directions of $X$ tensor to $E_2$ by the additivity
theorem $E_1 \otimes E_1 = E_2$ of Lurie. The $(-2)$-shift of
Theorem~\ref{thm:minus2-shift} is exactly the CY anomaly grading
that promotes the $E_2$-structure to a Poisson-$E_2$ structure in
the sense of Calaque--Pantev--To\"en--Vaqui\'e--Vezzosi; the bracket
$\{-,-\}$ lowers total degree by $2$.
\end{proof}

At $d=4$ the same procedure gives $E_0$-observables (just a pointed
space) with $0$-shifted pairing, a pure topological field theory
with no residual holomorphic content. This is why 4D hCS on CY$_2$
is the last CY-dimension at which $\Phi$ retains a nontrivial chiral
image: the shift equals $d-4$, and crossing $0$ collapses the data.

\section{$\kappach$ on K3 and $T^4$}

\begin{proposition}[CY$_2$ values]
\label{prop:kappa-ch-CY2}
For the classical observable algebra at the CY$_2$ fixed point,
\[
\kappach(\cA_X) \;=\; \sum_q (-1)^q h^{0,q}(X).
\]
Evaluation:
\begin{itemize}
\item $X = \mathrm{K3}$: $h^{0,\bullet} = (1,0,1)$, so
 $\kappach(\cA_{\mathrm{K3}}) = 1 - 0 + 1 = 2$.
\item $X = T^4$: $h^{0,\bullet} = (1,2,1)$, so
 $\kappach(\cA_{T^4}) = 1 - 2 + 1 = 0$.
\end{itemize}
Both agree with $\kappacat(X) = \chi(\cO_X)$: at $d=2$ the
Hodge-supertrace $\kappach$ coincides with $\kappacat$ because
$\kappacat$ is honestly the Euler characteristic of the structure
sheaf (canonically $2$ for K3, $0$ for $T^4$), and the alternating
sum $\sum_q (-1)^q h^{0,q}$ is the Hirzebruch--Riemann--Roch formula
for $\chi(\cO_X)$.
\end{proposition}

\begin{proof}
The first identification is the compact-CY supertrace identification
$\kappach(\cA_X) = \sum_q (-1)^q h^{0,q}(X)$ of the $E_n$-hierarchy
part of Vol~III. The Hodge numbers are classical: K3 has
$(1,0,1)$ since $h^{1,0}(K3)=0$ and $h^{2,0}(K3)=1$; $T^4$ has
$(1,2,1)$ since $H^1(T^4)$ is $4$-dimensional and splits as
$h^{1,0}+h^{0,1} = 2+2$, so $h^{0,1}(T^4) = 2$ and $h^{0,2}(T^4) =
\binom{2}{2} = 1$. See
\texttt{chapters/examples/cy\_d\_kappa\_stratification.tex}
lines $730$--$740$ for the tabulated entries.
\end{proof}

The T$^4$ cancellation is not an accident. Supertrace on an abelian
CY$_2$ sees the free $H^1$-space twice with opposite signs; the
bielliptic surface cancels for the same reason at
$(1,1,0)$, returning $0$. Only the K3 column $(1,0,1)$ survives
as a nontrivial CY$_2$ $\kappach$, a structural reason the K3 Yangian
and the Gritsenko $\Delta_5$ Borcherds weight are what they are.

\section{Class M and the CY-B $E_2$-chiral stratum}

\begin{proposition}[Class M match]
\label{prop:class-M-match}
The CY-B stratification of $\Phi$ (cf.
\texttt{chapters/theory/e2\_chiral\_algebras.tex} $\S398$) asserts
$\Phi_2(\cC)$ is natively $E_2$-chiral at $d=2$. Theorem~\ref{thm:E2-obs}
furnishes the $E_2$-structure on the observable side, and
Theorem~\ref{thm:minus2-shift} supplies the compatible $(-2)$-shifted
symplectic datum that gives $\Phi_2(\cC)$ its Poisson-$E_2$
bracket. The two statements are the BV and the $\infty$-categorical
lanes of the same theorem.
\end{proposition}

At $d=2$ the CY-B ambient category is Class M: 2-dimensional
Calabi--Yau categories whose image under $\Phi$ carries an
$E_2$-chiral algebra structure. For $X = \mathrm{K3}$,
$\Phi_2(\mathrm{Coh}(\mathrm{K3}))$ is the
Schiffmann--Vasserot K3 cohomological Hall algebra;
Theorem~\ref{thm:E2-obs} says its factorisation-algebra model from
4D hCS is $E_2$, and Theorem~\ref{thm:minus2-shift} says the
associated bracket is the $(-2)$-shifted Poisson bracket of PTVV.
The Class M $E_3$-bar dimension $6^g$ on a genus-$g$ Riemann surface
(cf. essential constants) is precisely the $E_2$-observable
dimension one level up, matching Lurie additivity $E_2\otimes E_1 = E_3$
in the factorisation-on-Riemann-surfaces direction.

\section{Shift = $d - 4$ across CY-A/B/C/D}

The table below compiles the $d$-shift law:
\[
\begin{array}{c|ccc}
d & \text{hCS dim} & \text{PTVV shift} & \text{Obs-structure} \\
\hline
2 & 4 & -2 & E_2 \\
3 & 6 & -1 & E_1 \\
4 & 8 & 0 & E_0
\end{array}
\]
Shift and $E_n$-index sum to $d-2$, which is the holomorphic
dimension of the disc directions after passing to Dolbeault
cohomology. At $d=4$ the theory degenerates: $0$-shifted means no
holomorphic residue, $E_0$ means no nontrivial algebraic structure.
4D hCS on CY$_2$ is the last nontrivial stratum; $d=5$ requires
$E_{-1}$, forbidden, hence the stop. Costello's CY$_2$-anchored
construction \cite{Costello2016} is recovered as the $d=2$ entry,
and the $(-2)$-shift is what Wave~14 I4's dim-shift rule predicts.

\begin{thebibliography}{9}
\bibitem{PTVV2013} T.~Pantev, B.~To\"en, M.~Vaqui\'e, G.~Vezzosi,
  \emph{Shifted symplectic structures}, Publ.~IHES 117 (2013), 271--328.
\bibitem{CPTVV2017} D.~Calaque, T.~Pantev, B.~To\"en, M.~Vaqui\'e,
  G.~Vezzosi, \emph{Shifted Poisson structures and deformation
  quantization}, J.~Topology 10 (2017), 483--584.
\bibitem{Costello2016} K.~Costello, \emph{M-theory in the $\Omega$-background
  and 5-dimensional non-commutative gauge theory}, arXiv:1610.04144 (2016).
\bibitem{CG2021} K.~Costello, O.~Gwilliam, \emph{Factorization Algebras in
  Quantum Field Theory, Volume II}, Cambridge University Press (2021).
\end{thebibliography}

\end{document}
